% Sección 8: Parte 5, opción B (sistema de recomendación e-commerce).
% Fragmento para main.tex: va con % Sección 8: Parte 5, opción B (sistema de recomendación e-commerce).
% Fragmento para main.tex: va con % Sección 8: Parte 5, opción B (sistema de recomendación e-commerce).
% Fragmento para main.tex: va con % Sección 8: Parte 5, opción B (sistema de recomendación e-commerce).
% Fragmento para main.tex: va con \input{seccion8_recomendador} justo después de la sección
% "Parte 5: Aplicación de inteligencia artificial" y antes de \begin{thebibliography}.
% Usa los macros del preámbulo de main.tex (\imagen, estilos TikZ caja/flecha, \citep).
% Las referencias nuevas están comentadas al final de este archivo: copiarlas a thebibliography.

\section{Parte 5: Sistema de recomendación e-commerce}
Además del RAG se implementó la opción B del enunciado: una tienda en línea de videojuegos que, para cada producto, recomienda un Top 10 combinando cinco señales \citep{burke}. Cada señal se resuelve con una estructura distinta del minigestor, así que la aplicación integra los cuatro tipos de datos del proyecto: relacional (órdenes, categoría y precio), espacial (ubicación de las tiendas), texto (descripciones) y multimedia (imágenes de los productos).

\begin{table}[H]
\centering\small
\caption{Señales del recomendador y estructura del minigestor que usa cada una}
\begin{tabularx}{\linewidth}{lYl}
\toprule
Señal & Pregunta que responde & Estructura \\
\midrule
Co-compra & ¿Qué compraron también quienes compraron este producto? & B+, self-join, \texttt{GROUP BY} \\
Texto & ¿Qué productos tienen una descripción parecida? & Índice invertido SPIMI, BM25 \\
Visual & ¿Qué productos se ven parecidos? & SIFT + BoVW, HNSW coseno \\
Metadata & ¿Qué productos de la misma categoría y precio están mejor valorados? & B+ sobre categoría \\
Espacial & ¿Qué productos hay en tiendas cerca de mí? & R-Tree: radio, k-NN y polígono \\
\bottomrule
\end{tabularx}
\end{table}

\subsection{Arquitectura}
La aplicación vive en la carpeta \texttt{ecommerce/} del repositorio y es independiente del motor. Se comunica con él solo mediante \texttt{POST /api/query}, igual que el RAG, y no modifica el código del minigestor.

\begin{figure}[H]
\centering
\resizebox{\linewidth}{!}{%
\begin{tikzpicture}[node distance=12mm and 14mm]
  \node[caja, minimum width=3cm] (nav) {\textbf{Tienda web}\\HTML + JS + Leaflet};
  \node[caja, right=of nav, minimum width=3.4cm] (srv) {\textbf{Servidor de la tienda}\\FastAPI, puerto 8300};
  \node[caja, right=of srv, minimum width=3.4cm] (rec) {\textbf{Recomendador}\\5 señales $\cdot$ fusión};
  \node[caja, below=of rec, minimum width=3.4cm] (vis) {\textbf{Vectores de imagen}\\SIFT k=4096 $\cdot$ CLIP\\local, CPU};
  \node[caja, below=of srv, minimum width=3.4cm, fill=gray!10] (motor) {\textbf{Minigestor BD2}\\puerto 8000\\B+ $\cdot$ SPIMI $\cdot$ HNSW $\cdot$ R-Tree};
  \draw[flecha] (nav) -- node[above, font=\scriptsize]{producto o foto} (srv);
  \draw[flecha] (srv) -- (rec);
  \draw[flecha] (rec) -- node[right, font=\scriptsize]{imagen} (vis);
  \draw[flecha] (rec) -- node[above left, font=\scriptsize]{SQL por HTTP} (motor);
  \draw[flecha, dashed] (rec.north) -- ++(0,0.6) -| node[above, pos=0.25, font=\scriptsize]{Top 10 con el aporte de cada señal} (nav.north);
\end{tikzpicture}}
\caption{Arquitectura del recomendador: la aplicación solo usa la API SQL del minigestor}\label{fig:rec}
\end{figure}

\begin{itemize}
  \item \texttt{recomendador.py} tiene una función por señal y la fusión. No guarda estado: cada recomendación consulta al motor.
  \item \texttt{cliente\_bd.py} es la única puerta hacia el motor. Arma los literales SQL de forma segura y anota el SQL, el plan y los tiempos de cada consulta, que la interfaz muestra.
  \item \texttt{visuales.py} calcula en la aplicación los vectores de los dos modelos de imagen que el motor no tiene cargados (sección~\ref{sec:rec-senales}).
\end{itemize}

\subsection{Dataset}
\begin{itemize}
  \item Se usó Amazon Reviews 2023 \citep{amazon2023}, categoría \texttt{Video\_Games}. Las interacciones salen del archivo \emph{5-core} (usuarios y productos con al menos 5 reseñas) y se tomaron los cerca de 4\,500 productos más reseñados como candidatos.
  \item El catálogo final tiene 1\,500 productos con imagen, precio (puede faltar), rating, marca y categoría. La categoría usa el segundo y tercer nivel del árbol de Amazon, por ejemplo \emph{Nintendo Switch > Games}, así que queda separada por plataforma.
  \item La descripción une título, marca, categoría, características y descripción del producto en un solo texto.
  \item Cada reseña se toma como una compra. Se generó una orden por par (usuario, producto), solo para usuarios con al menos 2 productos del catálogo: resultaron 24\,999 órdenes.
  \item Las 8 tiendas son puntos reales de Lima (Miraflores, San Isidro, Surco, San Borja, Barranco, Jesús María, San Miguel y Callao). Cada producto se asigna a una tienda de forma determinista con \texttt{crc32(codigo) mod 8}.
\end{itemize}

\subsection{Modelo de datos}
\begin{lstlisting}[language=SQL]
CREATE TABLE sucursales (id INT PRIMARY KEY, nombre VARCHAR(100),
  distrito VARCHAR(60), ubicacion POINT) USING BPLUS;
CREATE TABLE catalogo (id INT PRIMARY KEY, codigo VARCHAR(20),
  nombre VARCHAR(200), marca VARCHAR(80), categoria VARCHAR(80),
  precio FLOAT, rating FLOAT, sucursal_id INT, imagen VARCHAR(255),
  descripcion TEXT, emb VECTOR(256)) USING BPLUS;
CREATE TABLE ordenes (id INT PRIMARY KEY, usuario_id INT, item_id INT,
  fecha DATE) USING BPLUS;
CREATE INDEX idx_cat_categoria ON catalogo (categoria) USING BPLUS;
CREATE INDEX idx_cat_sucursal ON catalogo (sucursal_id) USING BPLUS;
CREATE INDEX ft_catalogo ON catalogo (descripcion) USING FULLTEXT;
CREATE INDEX hnsw_catalogo ON catalogo (emb) USING HNSW WITH (METRIC=cosine);
CREATE INDEX idx_ord_usuario ON ordenes (usuario_id) USING BPLUS;
CREATE INDEX idx_ord_item ON ordenes (item_id) USING BPLUS;
\end{lstlisting}
\begin{itemize}
  \item La columna \texttt{ubicacion} es \texttt{POINT}, así que el motor le crea su R-Tree automáticamente (parte 2).
  \item El vector \texttt{emb} lo calcula el propio motor a partir de la ruta de la imagen, con el pipeline SIFT de la parte 4.
  \item Los otros dos modelos de imagen viven en tablas propias, \texttt{vis\_sift4096} con \texttt{VECTOR(4096)} y \texttt{vis\_clip} con \texttt{VECTOR(512)}, cada una con su índice HNSW. Son opcionales.
  \item La ingesta inserta por lotes de 40 filas, crea los índices al final y corre \texttt{ANALYZE} en cada tabla para que el planificador tenga estadísticas.
\end{itemize}

\subsection{Señales de recomendación}\label{sec:rec-senales}
Todas las señales reciben los mismos filtros de categoría, precio y tiendas permitidas como condiciones extra del \texttt{WHERE}, excluyen al producto consultado y devuelven hasta 30 candidatos.

\paragraph{Co-compra.} Es filtrado colaborativo de producto a producto \citep{linden}: \emph{quienes compraron X también compraron Y}.
\begin{lstlisting}[language=SQL]
SELECT o2.item_id, COUNT(*) FROM ordenes AS o1
JOIN ordenes AS o2 ON o1.usuario_id = o2.usuario_id
WHERE o1.item_id = 42 AND o2.item_id != 42
GROUP BY o2.item_id ORDER BY COUNT(*) DESC LIMIT 30;
\end{lstlisting}
El plan usa el B+ de \texttt{ordenes.item\_id} para hallar a los compradores, un hash join externo, \texttt{GROUP BY} por hash y un top-N con heap que solo guarda los 30 mejores. Es la señal más fuerte y la más cara, unos 790~ms. En un producto sin historial (\emph{cold-start}) devuelve vacío.

\paragraph{Visual.} Compara la imagen del producto con las del catálogo, o una foto que sube el usuario. Hay tres modelos:
\begin{itemize}
  \item SIFT + BoVW con 256 palabras, el pipeline de la parte 4 \citep{lowe,sivic}: hasta 500 descriptores SIFT por imagen, vocabulario entrenado con el K-Means propio sobre 100\,000 descriptores del catálogo, histograma con TF-IDF e índice HNSW con coseno \citep{hnsw}. El motor extrae los descriptores de la consulta a partir de la ruta (unos 130~ms).
  \item SIFT + BoVW con 4\,096 palabras y TF sublineal. Como el motor carga un solo vocabulario a la vez, la aplicación calcula el vector reutilizando los módulos SIFT y TF-IDF del motor y lo busca con \texttt{SIMILAR\_TO(VECTOR(\ldots), k=30)}.
  \item CLIP ViT-B/32 \citep{clip}, una red preentrenada que da vectores de 512 dimensiones. Está solo para comparar contra SIFT.
\end{itemize}
Se agregaron los dos últimos porque con 256 palabras la búsqueda por foto fallaba en casos claros. Con la carátula de \emph{Pokémon Shield}, la de \emph{Pokémon Sword} quedaba en el puesto 273 de 1\,500: con tan pocas palabras todas las carátulas se parecen. Con la foto de un control DualShock negro, el blanco quedaba en el puesto 1\,352, porque SIFT describe la dirección del cambio de brillo y un borde negro sobre blanco produce descriptores casi opuestos a uno blanco sobre blanco. CLIP lo encuentra en el primer puesto.

\paragraph{Texto.} \texttt{MATCH(descripcion, \ldots) USING BM25 LIMIT 30} sobre el índice SPIMI de la parte 3 \citep{robertson}, usando como consulta las primeras 25 palabras del nombre, la marca, la categoría y la descripción del producto. El recorte mantiene el costo estable, en unos 30~ms.

\paragraph{Metadata.} Los productos mejor valorados de la misma categoría con un precio entre 0.5 y 1.8 veces el del producto consultado, completando con la misma categoría sin rango de precio si faltan candidatos. Usa el B+ de \texttt{categoria} y un top-N por rating (unos 31~ms).

\paragraph{Espacial.} La ubicación del usuario, elegida con un clic en el mapa, se usa de dos formas:
\begin{itemize}
  \item Como filtro: \texttt{distancia(ubicacion, POINT(\ldots)) < radio} o \texttt{intersecta(ubicacion, POLYGON(\ldots))} sobre el R-Tree \citep{guttman} devuelve las tiendas permitidas, que pasan a las demás señales como \texttt{sucursal\_id IN (\ldots)}. La distancia es Haversine \citep{sinnott}.
  \item Como señal: las 3 tiendas más cercanas (k-NN del R-Tree) dan a cada producto suyo un puntaje $1 / (1 + d / 5000\,\text{m})$. En el híbrido solo puntúa candidatos de las otras señales, para acercar el Top 10 sin agregar productos irrelevantes.
\end{itemize}

\subsection{Fusión híbrida}
Cada señal entrega sus candidatos con un puntaje en su propia escala (conteos, coseno, BM25 o rating). Se combinan con una de dos fórmulas, que el usuario elige:
\begin{itemize}
  \item Suma ponderada de puntajes normalizados, la opción por defecto:
  \[
  \text{puntaje}(p) = \sum_{m} w_m \, \frac{s_m(p)}{\max_q s_m(q)}
  \]
  La co-compra usa $\log(1 + \text{conteo})$ antes de normalizar, para que un producto con muchas co-compras no aplaste al resto. La señal espacial ya está entre 0 y 1 y no se normaliza.
  \item \emph{Reciprocal Rank Fusion} ponderada \citep{rrf}, la misma del RAG pero con un peso por señal:
  \[
  \text{RRF}(p) = \sum_{m} \frac{w_m}{60 + \text{rango}_m(p)}
  \]
\end{itemize}
Los pesos de las señales activas se renormalizan para sumar 1, y una señal con peso 0 no se consulta. Los empates se resuelven por mayor rating y luego por menor \texttt{id}. Los pesos por defecto se calibraron con los usuarios de validación (sección~\ref{sec:rec-eval}): co-compra 0.80, texto 0.15, visual 0.05 y metadata 0, más 0.10 de cercanía cuando hay ubicación. Cada recomendación devuelve cuánto aportó cada señal a su puntaje.

\subsection{Interfaz}
\begin{itemize}
  \item Catálogo paginado con filtros por categoría y precio, búsqueda por texto (BM25 o TF-IDF) y búsqueda por foto. En la búsqueda por foto, la imagen subida queda fija sobre los resultados y cada producto muestra su similitud coseno con ella.
  \item Al elegir un producto, un panel lateral muestra su ficha y el Top 10. Bajo cada recomendación, una franja de colores indica cuánto aportó cada señal.
  \item Un panel de preferencias permite cambiar la estrategia (híbrido o una señal sola), los pesos, la fórmula de fusión y el modelo de imagen sin cerrar el producto.
  \item Un mapa de Leaflet fija la ubicación del usuario, el radio y la zona, y marca las tiendas con recomendaciones y las más cercanas.
  \item Cada vista despliega el SQL enviado al motor, el plan ejecutado y los tiempos.
\end{itemize}
\imagen{imagenes/ecommerce/recomendacion.png}{Producto elegido y su Top 10 híbrido, con la franja de aporte de cada señal}{fig:rec-detalle}
\imagen{imagenes/ecommerce/busqueda_foto.png}{Búsqueda por foto: la imagen subida y los productos más parecidos con su similitud}{fig:rec-foto}

\subsection{Evaluación}\label{sec:rec-eval}
\subsubsection{Metodología}
\begin{itemize}
  \item Los usuarios con al menos 2 productos se repartieron al azar en entrenamiento (60\,\%), validación (20\,\%) y prueba (20\,\%), con una semilla fija.
  \item Por cada usuario de prueba (100 por corrida) se elige al azar uno de sus productos como consulta. Los objetivos son sus demás productos. La co-compra excluye las órdenes de los usuarios de prueba, para no usar la respuesta.
  \item El escenario \emph{cold-start} repite los mismos casos sin ninguna co-compra, como un producto nuevo.
  \item Los pesos del híbrido se eligieron solo con los 301 usuarios de validación, probando las 1\,771 combinaciones con paso 0.05 que suman 1. Se eligió la de mayor NDCG@10 promedio entre el escenario normal y el \emph{cold-start}. Optimizar solo el normal dejaba en 0 todas las señales de contenido.
  \item Métricas: Hit@10 (casos con al menos un objetivo en el Top 10), Recall@10, NDCG@10 \citep{ndcg}, MRR y Prec.~Cat@10, la fracción del Top 10 en la misma categoría que el producto consultado.
\end{itemize}

\subsubsection{Resultados}
\begin{table}[H]
\centering\small
\caption{Recomendación por señal y fusión, 100 usuarios de prueba}
\begin{tabular}{llccccc}
\toprule
Escenario & Modo & Hit@10 & Recall@10 & NDCG@10 & MRR & Prec. Cat@10 \\
\midrule
Normal & Co-compra & 0.60 & 0.083 & 0.134 & 0.324 & 0.304 \\
 & Visual (SIFT k=256) & 0.20 & 0.019 & 0.028 & 0.073 & 0.219 \\
 & Texto & 0.34 & 0.037 & 0.063 & 0.163 & 0.439 \\
 & Metadata & 0.15 & 0.014 & 0.017 & 0.035 & \textbf{1.000} \\
 & Híbrido RRF & \textbf{0.63} & 0.085 & \textbf{0.137} & \textbf{0.343} & 0.321 \\
 & Híbrido score & 0.61 & \textbf{0.086} & 0.135 & 0.323 & 0.323 \\
\midrule
Cold-start & Texto & \textbf{0.34} & \textbf{0.037} & \textbf{0.063} & \textbf{0.163} & 0.439 \\
 & Híbrido RRF & \textbf{0.34} & \textbf{0.037} & 0.060 & 0.145 & 0.440 \\
 & Híbrido score & \textbf{0.34} & \textbf{0.037} & 0.062 & 0.156 & 0.432 \\
\bottomrule
\end{tabular}
\end{table}
En \emph{cold-start} la co-compra no recomienda nada y las demás señales solas rinden igual que en el escenario normal. Con otra semilla de casos los resultados se mantienen: co-compra 0.61 de Hit@10 e híbrido score 0.63.

\begin{table}[H]
\centering\small
\caption{Modelos de imagen como señal sola, mismos casos}
\begin{tabular}{lccccc}
\toprule
Modelo & Hit@10 & Recall@10 & NDCG@10 & MRR & Prec. Cat@10 \\
\midrule
SIFT k=256 & 0.20 & 0.019 & 0.028 & 0.073 & 0.219 \\
SIFT k=4096 & 0.27 & 0.027 & 0.040 & 0.109 & \textbf{0.373} \\
CLIP & \textbf{0.38} & \textbf{0.044} & \textbf{0.067} & \textbf{0.171} & 0.288 \\
\bottomrule
\end{tabular}
\end{table}

\begin{table}[H]
\centering\small
\caption{Peso de la cercanía en el híbrido score, con el usuario en uno de 8 puntos de Lima}
\begin{tabular}{cccccc}
\toprule
Peso espacial & Hit@10 & NDCG@10 & MRR & Prec. Cat@10 & Distancia media del Top 10 \\
\midrule
0.00 & 0.61 & \textbf{0.135} & \textbf{0.323} & 0.323 & 7.86 km \\
0.10 & \textbf{0.62} & 0.130 & 0.297 & \textbf{0.328} & 6.77 km \\
0.25 & \textbf{0.62} & 0.121 & 0.279 & 0.318 & 5.63 km \\
0.50 & 0.58 & 0.107 & 0.247 & 0.295 & 4.68 km \\
\bottomrule
\end{tabular}
\end{table}

\begin{table}[H]
\centering\small
\caption{Tiempo en el motor de una recomendación híbrida con ubicación}
\begin{tabular}{llr}
\toprule
Consulta & Plan & Tiempo \\
\midrule
Co-compra (self-join) & B+ \texttt{item\_id}, hash join externo, top-N & $\sim$792 ms \\
Datos de los candidatos de co-compra & Sequential Scan & $\sim$127 ms \\
Visual & HNSW coseno & $\sim$133 ms \\
Texto & SPIMI BM25 & $\sim$30 ms \\
Metadata & B+ \texttt{categoria}, top-N & $\sim$31 ms \\
Radio y k-NN & R-Tree & $\sim$1 ms cada una \\
\bottomrule
\end{tabular}
\end{table}

\subsubsection{Análisis}
\begin{itemize}
  \item Con historial de compras, el híbrido iguala a la co-compra sola (Hit@10 entre 0.61 y 0.63 frente a 0.60 y 0.61) y sube algo la coherencia de categoría. Con 100 casos, diferencias de 0.02 están dentro del ruido.
  \item En \emph{cold-start} el híbrido rinde como el texto solo, la mejor señal en ese escenario. Con visual 0.05 y metadata 0, al faltar la co-compra casi solo queda el texto.
  \item Elegir los pesos es elegir una métrica. Unos pesos manuales con más metadata y visual (co-compra 0.55, texto 0.25, visual 0.08 y metadata 0.12) dan el mismo NDCG@10, pero más Hit@10 (0.66 en el escenario normal y 0.40 en \emph{cold-start}) y más coherencia de categoría (0.40 y 0.60, frente a 0.32 y 0.43). La calibración optimiza el NDCG sobre \emph{lo que el usuario también compró}, que favorece a la co-compra por diseño. Esos pesos manuales se eligieron antes de separar los usuarios al azar, así que su ventaja no es una medición limpia.
  \item Suma de puntajes y RRF quedan empatadas: cada una gana en una de las dos semillas, siempre dentro del ruido.
  \item Como señal sola, CLIP casi duplica el Hit@10 de SIFT con 256 palabras. SIFT con 4\,096 palabras tiene la mejor coherencia de categoría, porque reconoce el marco y el logo de la plataforma en la carátula. En el híbrido el modelo de imagen casi no cambia los resultados, porque su peso es 0.05; importa sobre todo en la búsqueda por foto.
  \item Un peso espacial de 0.10 acerca el Top 10 unos 1.1~km sin bajar el Hit@10. Desde 0.50 empeoran todas las métricas.
  \item La co-compra concentra cerca del 80\,\% del tiempo de una recomendación (unos 1\,000~ms en total). El R-Tree es prácticamente gratis.
\end{itemize}

\subsection{Alcance y limitaciones}
\begin{itemize}
  \item Un acierto significa que el usuario también compró el producto. La métrica mide solo de forma indirecta la similitud visual o de texto, y la coherencia de categoría la complementa.
  \item La co-compra no funciona con productos nuevos, y el self-join crece con la popularidad del producto. A mayor escala convendría precalcular los pares de co-compra.
  \item SIFT no encuentra el mismo objeto con otro color. CLIP sí, pero es un modelo preentrenado externo al minigestor.
  \item \texttt{WHERE id IN (\ldots)} sobre el catálogo se resuelve con un recorrido secuencial en vez del B+ de \texttt{id}. No se corrigió para no modificar el motor.
  \item La interfaz no registra compras, así que la co-compra no aprende en línea.
  \item La señal espacial en el híbrido solo reordena candidatos de las otras señales: no descubre productos nuevos.
\end{itemize}

% ==================== REFERENCIAS NUEVAS ====================
% Copiar dentro de \begin{thebibliography}{99} de main.tex, en orden alfabético.
% Las demás citas de esta sección (rrf, lowe, sivic, hnsw, robertson, guttman, sinnott) ya existen.
%
% \bibitem[Burke(2002)]{burke} Burke, R. (2002). Hybrid recommender systems: survey and experiments. \textit{User Modeling and User-Adapted Interaction}, 12(4), 331--370.
% \bibitem[Hou et~al.(2024)]{amazon2023} Hou, Y., Li, J., He, Z., Yan, A., Chen, X. y McAuley, J. (2024). Bridging language and items for retrieval and recommendation. \textit{arXiv:2403.03952}. Dataset: \url{https://amazon-reviews-2023.github.io/}
% \bibitem[Järvelin y Kekäläinen(2002)]{ndcg} Järvelin, K. y Kekäläinen, J. (2002). Cumulated gain-based evaluation of IR techniques. \textit{ACM Transactions on Information Systems}, 20(4), 422--446.
% \bibitem[Linden et~al.(2003)]{linden} Linden, G., Smith, B. y York, J. (2003). Amazon.com recommendations: item-to-item collaborative filtering. \textit{IEEE Internet Computing}, 7(1), 76--80.
% \bibitem[Radford et~al.(2021)]{clip} Radford, A., Kim, J. W., Hallacy, C., Ramesh, A., Goh, G., Agarwal, S., Sastry, G., Askell, A., Mishkin, P., Clark, J., Krueger, G. y Sutskever, I. (2021). Learning transferable visual models from natural language supervision. \textit{Proceedings of ICML}, PMLR 139, 8748--8763.
 justo después de la sección
% "Parte 5: Aplicación de inteligencia artificial" y antes de \begin{thebibliography}.
% Usa los macros del preámbulo de main.tex (\imagen, estilos TikZ caja/flecha, \citep).
% Las referencias nuevas están comentadas al final de este archivo: copiarlas a thebibliography.

\section{Parte 5: Sistema de recomendación e-commerce}
Además del RAG se implementó la opción B del enunciado: una tienda en línea de videojuegos que, para cada producto, recomienda un Top 10 combinando cinco señales \citep{burke}. Cada señal se resuelve con una estructura distinta del minigestor, así que la aplicación integra los cuatro tipos de datos del proyecto: relacional (órdenes, categoría y precio), espacial (ubicación de las tiendas), texto (descripciones) y multimedia (imágenes de los productos).

\begin{table}[H]
\centering\small
\caption{Señales del recomendador y estructura del minigestor que usa cada una}
\begin{tabularx}{\linewidth}{lYl}
\toprule
Señal & Pregunta que responde & Estructura \\
\midrule
Co-compra & ¿Qué compraron también quienes compraron este producto? & B+, self-join, \texttt{GROUP BY} \\
Texto & ¿Qué productos tienen una descripción parecida? & Índice invertido SPIMI, BM25 \\
Visual & ¿Qué productos se ven parecidos? & SIFT + BoVW, HNSW coseno \\
Metadata & ¿Qué productos de la misma categoría y precio están mejor valorados? & B+ sobre categoría \\
Espacial & ¿Qué productos hay en tiendas cerca de mí? & R-Tree: radio, k-NN y polígono \\
\bottomrule
\end{tabularx}
\end{table}

\subsection{Arquitectura}
La aplicación vive en la carpeta \texttt{ecommerce/} del repositorio y es independiente del motor. Se comunica con él solo mediante \texttt{POST /api/query}, igual que el RAG, y no modifica el código del minigestor.

\begin{figure}[H]
\centering
\resizebox{\linewidth}{!}{%
\begin{tikzpicture}[node distance=12mm and 14mm]
  \node[caja, minimum width=3cm] (nav) {\textbf{Tienda web}\\HTML + JS + Leaflet};
  \node[caja, right=of nav, minimum width=3.4cm] (srv) {\textbf{Servidor de la tienda}\\FastAPI, puerto 8300};
  \node[caja, right=of srv, minimum width=3.4cm] (rec) {\textbf{Recomendador}\\5 señales $\cdot$ fusión};
  \node[caja, below=of rec, minimum width=3.4cm] (vis) {\textbf{Vectores de imagen}\\SIFT k=4096 $\cdot$ CLIP\\local, CPU};
  \node[caja, below=of srv, minimum width=3.4cm, fill=gray!10] (motor) {\textbf{Minigestor BD2}\\puerto 8000\\B+ $\cdot$ SPIMI $\cdot$ HNSW $\cdot$ R-Tree};
  \draw[flecha] (nav) -- node[above, font=\scriptsize]{producto o foto} (srv);
  \draw[flecha] (srv) -- (rec);
  \draw[flecha] (rec) -- node[right, font=\scriptsize]{imagen} (vis);
  \draw[flecha] (rec) -- node[above left, font=\scriptsize]{SQL por HTTP} (motor);
  \draw[flecha, dashed] (rec.north) -- ++(0,0.6) -| node[above, pos=0.25, font=\scriptsize]{Top 10 con el aporte de cada señal} (nav.north);
\end{tikzpicture}}
\caption{Arquitectura del recomendador: la aplicación solo usa la API SQL del minigestor}\label{fig:rec}
\end{figure}

\begin{itemize}
  \item \texttt{recomendador.py} tiene una función por señal y la fusión. No guarda estado: cada recomendación consulta al motor.
  \item \texttt{cliente\_bd.py} es la única puerta hacia el motor. Arma los literales SQL de forma segura y anota el SQL, el plan y los tiempos de cada consulta, que la interfaz muestra.
  \item \texttt{visuales.py} calcula en la aplicación los vectores de los dos modelos de imagen que el motor no tiene cargados (sección~\ref{sec:rec-senales}).
\end{itemize}

\subsection{Dataset}
\begin{itemize}
  \item Se usó Amazon Reviews 2023 \citep{amazon2023}, categoría \texttt{Video\_Games}. Las interacciones salen del archivo \emph{5-core} (usuarios y productos con al menos 5 reseñas) y se tomaron los cerca de 4\,500 productos más reseñados como candidatos.
  \item El catálogo final tiene 1\,500 productos con imagen, precio (puede faltar), rating, marca y categoría. La categoría usa el segundo y tercer nivel del árbol de Amazon, por ejemplo \emph{Nintendo Switch > Games}, así que queda separada por plataforma.
  \item La descripción une título, marca, categoría, características y descripción del producto en un solo texto.
  \item Cada reseña se toma como una compra. Se generó una orden por par (usuario, producto), solo para usuarios con al menos 2 productos del catálogo: resultaron 24\,999 órdenes.
  \item Las 8 tiendas son puntos reales de Lima (Miraflores, San Isidro, Surco, San Borja, Barranco, Jesús María, San Miguel y Callao). Cada producto se asigna a una tienda de forma determinista con \texttt{crc32(codigo) mod 8}.
\end{itemize}

\subsection{Modelo de datos}
\begin{lstlisting}[language=SQL]
CREATE TABLE sucursales (id INT PRIMARY KEY, nombre VARCHAR(100),
  distrito VARCHAR(60), ubicacion POINT) USING BPLUS;
CREATE TABLE catalogo (id INT PRIMARY KEY, codigo VARCHAR(20),
  nombre VARCHAR(200), marca VARCHAR(80), categoria VARCHAR(80),
  precio FLOAT, rating FLOAT, sucursal_id INT, imagen VARCHAR(255),
  descripcion TEXT, emb VECTOR(256)) USING BPLUS;
CREATE TABLE ordenes (id INT PRIMARY KEY, usuario_id INT, item_id INT,
  fecha DATE) USING BPLUS;
CREATE INDEX idx_cat_categoria ON catalogo (categoria) USING BPLUS;
CREATE INDEX idx_cat_sucursal ON catalogo (sucursal_id) USING BPLUS;
CREATE INDEX ft_catalogo ON catalogo (descripcion) USING FULLTEXT;
CREATE INDEX hnsw_catalogo ON catalogo (emb) USING HNSW WITH (METRIC=cosine);
CREATE INDEX idx_ord_usuario ON ordenes (usuario_id) USING BPLUS;
CREATE INDEX idx_ord_item ON ordenes (item_id) USING BPLUS;
\end{lstlisting}
\begin{itemize}
  \item La columna \texttt{ubicacion} es \texttt{POINT}, así que el motor le crea su R-Tree automáticamente (parte 2).
  \item El vector \texttt{emb} lo calcula el propio motor a partir de la ruta de la imagen, con el pipeline SIFT de la parte 4.
  \item Los otros dos modelos de imagen viven en tablas propias, \texttt{vis\_sift4096} con \texttt{VECTOR(4096)} y \texttt{vis\_clip} con \texttt{VECTOR(512)}, cada una con su índice HNSW. Son opcionales.
  \item La ingesta inserta por lotes de 40 filas, crea los índices al final y corre \texttt{ANALYZE} en cada tabla para que el planificador tenga estadísticas.
\end{itemize}

\subsection{Señales de recomendación}\label{sec:rec-senales}
Todas las señales reciben los mismos filtros de categoría, precio y tiendas permitidas como condiciones extra del \texttt{WHERE}, excluyen al producto consultado y devuelven hasta 30 candidatos.

\paragraph{Co-compra.} Es filtrado colaborativo de producto a producto \citep{linden}: \emph{quienes compraron X también compraron Y}.
\begin{lstlisting}[language=SQL]
SELECT o2.item_id, COUNT(*) FROM ordenes AS o1
JOIN ordenes AS o2 ON o1.usuario_id = o2.usuario_id
WHERE o1.item_id = 42 AND o2.item_id != 42
GROUP BY o2.item_id ORDER BY COUNT(*) DESC LIMIT 30;
\end{lstlisting}
El plan usa el B+ de \texttt{ordenes.item\_id} para hallar a los compradores, un hash join externo, \texttt{GROUP BY} por hash y un top-N con heap que solo guarda los 30 mejores. Es la señal más fuerte y la más cara, unos 790~ms. En un producto sin historial (\emph{cold-start}) devuelve vacío.

\paragraph{Visual.} Compara la imagen del producto con las del catálogo, o una foto que sube el usuario. Hay tres modelos:
\begin{itemize}
  \item SIFT + BoVW con 256 palabras, el pipeline de la parte 4 \citep{lowe,sivic}: hasta 500 descriptores SIFT por imagen, vocabulario entrenado con el K-Means propio sobre 100\,000 descriptores del catálogo, histograma con TF-IDF e índice HNSW con coseno \citep{hnsw}. El motor extrae los descriptores de la consulta a partir de la ruta (unos 130~ms).
  \item SIFT + BoVW con 4\,096 palabras y TF sublineal. Como el motor carga un solo vocabulario a la vez, la aplicación calcula el vector reutilizando los módulos SIFT y TF-IDF del motor y lo busca con \texttt{SIMILAR\_TO(VECTOR(\ldots), k=30)}.
  \item CLIP ViT-B/32 \citep{clip}, una red preentrenada que da vectores de 512 dimensiones. Está solo para comparar contra SIFT.
\end{itemize}
Se agregaron los dos últimos porque con 256 palabras la búsqueda por foto fallaba en casos claros. Con la carátula de \emph{Pokémon Shield}, la de \emph{Pokémon Sword} quedaba en el puesto 273 de 1\,500: con tan pocas palabras todas las carátulas se parecen. Con la foto de un control DualShock negro, el blanco quedaba en el puesto 1\,352, porque SIFT describe la dirección del cambio de brillo y un borde negro sobre blanco produce descriptores casi opuestos a uno blanco sobre blanco. CLIP lo encuentra en el primer puesto.

\paragraph{Texto.} \texttt{MATCH(descripcion, \ldots) USING BM25 LIMIT 30} sobre el índice SPIMI de la parte 3 \citep{robertson}, usando como consulta las primeras 25 palabras del nombre, la marca, la categoría y la descripción del producto. El recorte mantiene el costo estable, en unos 30~ms.

\paragraph{Metadata.} Los productos mejor valorados de la misma categoría con un precio entre 0.5 y 1.8 veces el del producto consultado, completando con la misma categoría sin rango de precio si faltan candidatos. Usa el B+ de \texttt{categoria} y un top-N por rating (unos 31~ms).

\paragraph{Espacial.} La ubicación del usuario, elegida con un clic en el mapa, se usa de dos formas:
\begin{itemize}
  \item Como filtro: \texttt{distancia(ubicacion, POINT(\ldots)) < radio} o \texttt{intersecta(ubicacion, POLYGON(\ldots))} sobre el R-Tree \citep{guttman} devuelve las tiendas permitidas, que pasan a las demás señales como \texttt{sucursal\_id IN (\ldots)}. La distancia es Haversine \citep{sinnott}.
  \item Como señal: las 3 tiendas más cercanas (k-NN del R-Tree) dan a cada producto suyo un puntaje $1 / (1 + d / 5000\,\text{m})$. En el híbrido solo puntúa candidatos de las otras señales, para acercar el Top 10 sin agregar productos irrelevantes.
\end{itemize}

\subsection{Fusión híbrida}
Cada señal entrega sus candidatos con un puntaje en su propia escala (conteos, coseno, BM25 o rating). Se combinan con una de dos fórmulas, que el usuario elige:
\begin{itemize}
  \item Suma ponderada de puntajes normalizados, la opción por defecto:
  \[
  \text{puntaje}(p) = \sum_{m} w_m \, \frac{s_m(p)}{\max_q s_m(q)}
  \]
  La co-compra usa $\log(1 + \text{conteo})$ antes de normalizar, para que un producto con muchas co-compras no aplaste al resto. La señal espacial ya está entre 0 y 1 y no se normaliza.
  \item \emph{Reciprocal Rank Fusion} ponderada \citep{rrf}, la misma del RAG pero con un peso por señal:
  \[
  \text{RRF}(p) = \sum_{m} \frac{w_m}{60 + \text{rango}_m(p)}
  \]
\end{itemize}
Los pesos de las señales activas se renormalizan para sumar 1, y una señal con peso 0 no se consulta. Los empates se resuelven por mayor rating y luego por menor \texttt{id}. Los pesos por defecto se calibraron con los usuarios de validación (sección~\ref{sec:rec-eval}): co-compra 0.80, texto 0.15, visual 0.05 y metadata 0, más 0.10 de cercanía cuando hay ubicación. Cada recomendación devuelve cuánto aportó cada señal a su puntaje.

\subsection{Interfaz}
\begin{itemize}
  \item Catálogo paginado con filtros por categoría y precio, búsqueda por texto (BM25 o TF-IDF) y búsqueda por foto. En la búsqueda por foto, la imagen subida queda fija sobre los resultados y cada producto muestra su similitud coseno con ella.
  \item Al elegir un producto, un panel lateral muestra su ficha y el Top 10. Bajo cada recomendación, una franja de colores indica cuánto aportó cada señal.
  \item Un panel de preferencias permite cambiar la estrategia (híbrido o una señal sola), los pesos, la fórmula de fusión y el modelo de imagen sin cerrar el producto.
  \item Un mapa de Leaflet fija la ubicación del usuario, el radio y la zona, y marca las tiendas con recomendaciones y las más cercanas.
  \item Cada vista despliega el SQL enviado al motor, el plan ejecutado y los tiempos.
\end{itemize}
\imagen{imagenes/ecommerce/recomendacion.png}{Producto elegido y su Top 10 híbrido, con la franja de aporte de cada señal}{fig:rec-detalle}
\imagen{imagenes/ecommerce/busqueda_foto.png}{Búsqueda por foto: la imagen subida y los productos más parecidos con su similitud}{fig:rec-foto}

\subsection{Evaluación}\label{sec:rec-eval}
\subsubsection{Metodología}
\begin{itemize}
  \item Los usuarios con al menos 2 productos se repartieron al azar en entrenamiento (60\,\%), validación (20\,\%) y prueba (20\,\%), con una semilla fija.
  \item Por cada usuario de prueba (100 por corrida) se elige al azar uno de sus productos como consulta. Los objetivos son sus demás productos. La co-compra excluye las órdenes de los usuarios de prueba, para no usar la respuesta.
  \item El escenario \emph{cold-start} repite los mismos casos sin ninguna co-compra, como un producto nuevo.
  \item Los pesos del híbrido se eligieron solo con los 301 usuarios de validación, probando las 1\,771 combinaciones con paso 0.05 que suman 1. Se eligió la de mayor NDCG@10 promedio entre el escenario normal y el \emph{cold-start}. Optimizar solo el normal dejaba en 0 todas las señales de contenido.
  \item Métricas: Hit@10 (casos con al menos un objetivo en el Top 10), Recall@10, NDCG@10 \citep{ndcg}, MRR y Prec.~Cat@10, la fracción del Top 10 en la misma categoría que el producto consultado.
\end{itemize}

\subsubsection{Resultados}
\begin{table}[H]
\centering\small
\caption{Recomendación por señal y fusión, 100 usuarios de prueba}
\begin{tabular}{llccccc}
\toprule
Escenario & Modo & Hit@10 & Recall@10 & NDCG@10 & MRR & Prec. Cat@10 \\
\midrule
Normal & Co-compra & 0.60 & 0.083 & 0.134 & 0.324 & 0.304 \\
 & Visual (SIFT k=256) & 0.20 & 0.019 & 0.028 & 0.073 & 0.219 \\
 & Texto & 0.34 & 0.037 & 0.063 & 0.163 & 0.439 \\
 & Metadata & 0.15 & 0.014 & 0.017 & 0.035 & \textbf{1.000} \\
 & Híbrido RRF & \textbf{0.63} & 0.085 & \textbf{0.137} & \textbf{0.343} & 0.321 \\
 & Híbrido score & 0.61 & \textbf{0.086} & 0.135 & 0.323 & 0.323 \\
\midrule
Cold-start & Texto & \textbf{0.34} & \textbf{0.037} & \textbf{0.063} & \textbf{0.163} & 0.439 \\
 & Híbrido RRF & \textbf{0.34} & \textbf{0.037} & 0.060 & 0.145 & 0.440 \\
 & Híbrido score & \textbf{0.34} & \textbf{0.037} & 0.062 & 0.156 & 0.432 \\
\bottomrule
\end{tabular}
\end{table}
En \emph{cold-start} la co-compra no recomienda nada y las demás señales solas rinden igual que en el escenario normal. Con otra semilla de casos los resultados se mantienen: co-compra 0.61 de Hit@10 e híbrido score 0.63.

\begin{table}[H]
\centering\small
\caption{Modelos de imagen como señal sola, mismos casos}
\begin{tabular}{lccccc}
\toprule
Modelo & Hit@10 & Recall@10 & NDCG@10 & MRR & Prec. Cat@10 \\
\midrule
SIFT k=256 & 0.20 & 0.019 & 0.028 & 0.073 & 0.219 \\
SIFT k=4096 & 0.27 & 0.027 & 0.040 & 0.109 & \textbf{0.373} \\
CLIP & \textbf{0.38} & \textbf{0.044} & \textbf{0.067} & \textbf{0.171} & 0.288 \\
\bottomrule
\end{tabular}
\end{table}

\begin{table}[H]
\centering\small
\caption{Peso de la cercanía en el híbrido score, con el usuario en uno de 8 puntos de Lima}
\begin{tabular}{cccccc}
\toprule
Peso espacial & Hit@10 & NDCG@10 & MRR & Prec. Cat@10 & Distancia media del Top 10 \\
\midrule
0.00 & 0.61 & \textbf{0.135} & \textbf{0.323} & 0.323 & 7.86 km \\
0.10 & \textbf{0.62} & 0.130 & 0.297 & \textbf{0.328} & 6.77 km \\
0.25 & \textbf{0.62} & 0.121 & 0.279 & 0.318 & 5.63 km \\
0.50 & 0.58 & 0.107 & 0.247 & 0.295 & 4.68 km \\
\bottomrule
\end{tabular}
\end{table}

\begin{table}[H]
\centering\small
\caption{Tiempo en el motor de una recomendación híbrida con ubicación}
\begin{tabular}{llr}
\toprule
Consulta & Plan & Tiempo \\
\midrule
Co-compra (self-join) & B+ \texttt{item\_id}, hash join externo, top-N & $\sim$792 ms \\
Datos de los candidatos de co-compra & Sequential Scan & $\sim$127 ms \\
Visual & HNSW coseno & $\sim$133 ms \\
Texto & SPIMI BM25 & $\sim$30 ms \\
Metadata & B+ \texttt{categoria}, top-N & $\sim$31 ms \\
Radio y k-NN & R-Tree & $\sim$1 ms cada una \\
\bottomrule
\end{tabular}
\end{table}

\subsubsection{Análisis}
\begin{itemize}
  \item Con historial de compras, el híbrido iguala a la co-compra sola (Hit@10 entre 0.61 y 0.63 frente a 0.60 y 0.61) y sube algo la coherencia de categoría. Con 100 casos, diferencias de 0.02 están dentro del ruido.
  \item En \emph{cold-start} el híbrido rinde como el texto solo, la mejor señal en ese escenario. Con visual 0.05 y metadata 0, al faltar la co-compra casi solo queda el texto.
  \item Elegir los pesos es elegir una métrica. Unos pesos manuales con más metadata y visual (co-compra 0.55, texto 0.25, visual 0.08 y metadata 0.12) dan el mismo NDCG@10, pero más Hit@10 (0.66 en el escenario normal y 0.40 en \emph{cold-start}) y más coherencia de categoría (0.40 y 0.60, frente a 0.32 y 0.43). La calibración optimiza el NDCG sobre \emph{lo que el usuario también compró}, que favorece a la co-compra por diseño. Esos pesos manuales se eligieron antes de separar los usuarios al azar, así que su ventaja no es una medición limpia.
  \item Suma de puntajes y RRF quedan empatadas: cada una gana en una de las dos semillas, siempre dentro del ruido.
  \item Como señal sola, CLIP casi duplica el Hit@10 de SIFT con 256 palabras. SIFT con 4\,096 palabras tiene la mejor coherencia de categoría, porque reconoce el marco y el logo de la plataforma en la carátula. En el híbrido el modelo de imagen casi no cambia los resultados, porque su peso es 0.05; importa sobre todo en la búsqueda por foto.
  \item Un peso espacial de 0.10 acerca el Top 10 unos 1.1~km sin bajar el Hit@10. Desde 0.50 empeoran todas las métricas.
  \item La co-compra concentra cerca del 80\,\% del tiempo de una recomendación (unos 1\,000~ms en total). El R-Tree es prácticamente gratis.
\end{itemize}

\subsection{Alcance y limitaciones}
\begin{itemize}
  \item Un acierto significa que el usuario también compró el producto. La métrica mide solo de forma indirecta la similitud visual o de texto, y la coherencia de categoría la complementa.
  \item La co-compra no funciona con productos nuevos, y el self-join crece con la popularidad del producto. A mayor escala convendría precalcular los pares de co-compra.
  \item SIFT no encuentra el mismo objeto con otro color. CLIP sí, pero es un modelo preentrenado externo al minigestor.
  \item \texttt{WHERE id IN (\ldots)} sobre el catálogo se resuelve con un recorrido secuencial en vez del B+ de \texttt{id}. No se corrigió para no modificar el motor.
  \item La interfaz no registra compras, así que la co-compra no aprende en línea.
  \item La señal espacial en el híbrido solo reordena candidatos de las otras señales: no descubre productos nuevos.
\end{itemize}

% ==================== REFERENCIAS NUEVAS ====================
% Copiar dentro de \begin{thebibliography}{99} de main.tex, en orden alfabético.
% Las demás citas de esta sección (rrf, lowe, sivic, hnsw, robertson, guttman, sinnott) ya existen.
%
% \bibitem[Burke(2002)]{burke} Burke, R. (2002). Hybrid recommender systems: survey and experiments. \textit{User Modeling and User-Adapted Interaction}, 12(4), 331--370.
% \bibitem[Hou et~al.(2024)]{amazon2023} Hou, Y., Li, J., He, Z., Yan, A., Chen, X. y McAuley, J. (2024). Bridging language and items for retrieval and recommendation. \textit{arXiv:2403.03952}. Dataset: \url{https://amazon-reviews-2023.github.io/}
% \bibitem[Järvelin y Kekäläinen(2002)]{ndcg} Järvelin, K. y Kekäläinen, J. (2002). Cumulated gain-based evaluation of IR techniques. \textit{ACM Transactions on Information Systems}, 20(4), 422--446.
% \bibitem[Linden et~al.(2003)]{linden} Linden, G., Smith, B. y York, J. (2003). Amazon.com recommendations: item-to-item collaborative filtering. \textit{IEEE Internet Computing}, 7(1), 76--80.
% \bibitem[Radford et~al.(2021)]{clip} Radford, A., Kim, J. W., Hallacy, C., Ramesh, A., Goh, G., Agarwal, S., Sastry, G., Askell, A., Mishkin, P., Clark, J., Krueger, G. y Sutskever, I. (2021). Learning transferable visual models from natural language supervision. \textit{Proceedings of ICML}, PMLR 139, 8748--8763.
 justo después de la sección
% "Parte 5: Aplicación de inteligencia artificial" y antes de \begin{thebibliography}.
% Usa los macros del preámbulo de main.tex (\imagen, estilos TikZ caja/flecha, \citep).
% Las referencias nuevas están comentadas al final de este archivo: copiarlas a thebibliography.

\section{Parte 5: Sistema de recomendación e-commerce}
Además del RAG se implementó la opción B del enunciado: una tienda en línea de videojuegos que, para cada producto, recomienda un Top 10 combinando cinco señales \citep{burke}. Cada señal se resuelve con una estructura distinta del minigestor, así que la aplicación integra los cuatro tipos de datos del proyecto: relacional (órdenes, categoría y precio), espacial (ubicación de las tiendas), texto (descripciones) y multimedia (imágenes de los productos).

\begin{table}[H]
\centering\small
\caption{Señales del recomendador y estructura del minigestor que usa cada una}
\begin{tabularx}{\linewidth}{lYl}
\toprule
Señal & Pregunta que responde & Estructura \\
\midrule
Co-compra & ¿Qué compraron también quienes compraron este producto? & B+, self-join, \texttt{GROUP BY} \\
Texto & ¿Qué productos tienen una descripción parecida? & Índice invertido SPIMI, BM25 \\
Visual & ¿Qué productos se ven parecidos? & SIFT + BoVW, HNSW coseno \\
Metadata & ¿Qué productos de la misma categoría y precio están mejor valorados? & B+ sobre categoría \\
Espacial & ¿Qué productos hay en tiendas cerca de mí? & R-Tree: radio, k-NN y polígono \\
\bottomrule
\end{tabularx}
\end{table}

\subsection{Arquitectura}
La aplicación vive en la carpeta \texttt{ecommerce/} del repositorio y es independiente del motor. Se comunica con él solo mediante \texttt{POST /api/query}, igual que el RAG, y no modifica el código del minigestor.

\begin{figure}[H]
\centering
\resizebox{\linewidth}{!}{%
\begin{tikzpicture}[node distance=12mm and 14mm]
  \node[caja, minimum width=3cm] (nav) {\textbf{Tienda web}\\HTML + JS + Leaflet};
  \node[caja, right=of nav, minimum width=3.4cm] (srv) {\textbf{Servidor de la tienda}\\FastAPI, puerto 8300};
  \node[caja, right=of srv, minimum width=3.4cm] (rec) {\textbf{Recomendador}\\5 señales $\cdot$ fusión};
  \node[caja, below=of rec, minimum width=3.4cm] (vis) {\textbf{Vectores de imagen}\\SIFT k=4096 $\cdot$ CLIP\\local, CPU};
  \node[caja, below=of srv, minimum width=3.4cm, fill=gray!10] (motor) {\textbf{Minigestor BD2}\\puerto 8000\\B+ $\cdot$ SPIMI $\cdot$ HNSW $\cdot$ R-Tree};
  \draw[flecha] (nav) -- node[above, font=\scriptsize]{producto o foto} (srv);
  \draw[flecha] (srv) -- (rec);
  \draw[flecha] (rec) -- node[right, font=\scriptsize]{imagen} (vis);
  \draw[flecha] (rec) -- node[above left, font=\scriptsize]{SQL por HTTP} (motor);
  \draw[flecha, dashed] (rec.north) -- ++(0,0.6) -| node[above, pos=0.25, font=\scriptsize]{Top 10 con el aporte de cada señal} (nav.north);
\end{tikzpicture}}
\caption{Arquitectura del recomendador: la aplicación solo usa la API SQL del minigestor}\label{fig:rec}
\end{figure}

\begin{itemize}
  \item \texttt{recomendador.py} tiene una función por señal y la fusión. No guarda estado: cada recomendación consulta al motor.
  \item \texttt{cliente\_bd.py} es la única puerta hacia el motor. Arma los literales SQL de forma segura y anota el SQL, el plan y los tiempos de cada consulta, que la interfaz muestra.
  \item \texttt{visuales.py} calcula en la aplicación los vectores de los dos modelos de imagen que el motor no tiene cargados (sección~\ref{sec:rec-senales}).
\end{itemize}

\subsection{Dataset}
\begin{itemize}
  \item Se usó Amazon Reviews 2023 \citep{amazon2023}, categoría \texttt{Video\_Games}. Las interacciones salen del archivo \emph{5-core} (usuarios y productos con al menos 5 reseñas) y se tomaron los cerca de 4\,500 productos más reseñados como candidatos.
  \item El catálogo final tiene 1\,500 productos con imagen, precio (puede faltar), rating, marca y categoría. La categoría usa el segundo y tercer nivel del árbol de Amazon, por ejemplo \emph{Nintendo Switch > Games}, así que queda separada por plataforma.
  \item La descripción une título, marca, categoría, características y descripción del producto en un solo texto.
  \item Cada reseña se toma como una compra. Se generó una orden por par (usuario, producto), solo para usuarios con al menos 2 productos del catálogo: resultaron 24\,999 órdenes.
  \item Las 8 tiendas son puntos reales de Lima (Miraflores, San Isidro, Surco, San Borja, Barranco, Jesús María, San Miguel y Callao). Cada producto se asigna a una tienda de forma determinista con \texttt{crc32(codigo) mod 8}.
\end{itemize}

\subsection{Modelo de datos}
\begin{lstlisting}[language=SQL]
CREATE TABLE sucursales (id INT PRIMARY KEY, nombre VARCHAR(100),
  distrito VARCHAR(60), ubicacion POINT) USING BPLUS;
CREATE TABLE catalogo (id INT PRIMARY KEY, codigo VARCHAR(20),
  nombre VARCHAR(200), marca VARCHAR(80), categoria VARCHAR(80),
  precio FLOAT, rating FLOAT, sucursal_id INT, imagen VARCHAR(255),
  descripcion TEXT, emb VECTOR(256)) USING BPLUS;
CREATE TABLE ordenes (id INT PRIMARY KEY, usuario_id INT, item_id INT,
  fecha DATE) USING BPLUS;
CREATE INDEX idx_cat_categoria ON catalogo (categoria) USING BPLUS;
CREATE INDEX idx_cat_sucursal ON catalogo (sucursal_id) USING BPLUS;
CREATE INDEX ft_catalogo ON catalogo (descripcion) USING FULLTEXT;
CREATE INDEX hnsw_catalogo ON catalogo (emb) USING HNSW WITH (METRIC=cosine);
CREATE INDEX idx_ord_usuario ON ordenes (usuario_id) USING BPLUS;
CREATE INDEX idx_ord_item ON ordenes (item_id) USING BPLUS;
\end{lstlisting}
\begin{itemize}
  \item La columna \texttt{ubicacion} es \texttt{POINT}, así que el motor le crea su R-Tree automáticamente (parte 2).
  \item El vector \texttt{emb} lo calcula el propio motor a partir de la ruta de la imagen, con el pipeline SIFT de la parte 4.
  \item Los otros dos modelos de imagen viven en tablas propias, \texttt{vis\_sift4096} con \texttt{VECTOR(4096)} y \texttt{vis\_clip} con \texttt{VECTOR(512)}, cada una con su índice HNSW. Son opcionales.
  \item La ingesta inserta por lotes de 40 filas, crea los índices al final y corre \texttt{ANALYZE} en cada tabla para que el planificador tenga estadísticas.
\end{itemize}

\subsection{Señales de recomendación}\label{sec:rec-senales}
Todas las señales reciben los mismos filtros de categoría, precio y tiendas permitidas como condiciones extra del \texttt{WHERE}, excluyen al producto consultado y devuelven hasta 30 candidatos.

\paragraph{Co-compra.} Es filtrado colaborativo de producto a producto \citep{linden}: \emph{quienes compraron X también compraron Y}.
\begin{lstlisting}[language=SQL]
SELECT o2.item_id, COUNT(*) FROM ordenes AS o1
JOIN ordenes AS o2 ON o1.usuario_id = o2.usuario_id
WHERE o1.item_id = 42 AND o2.item_id != 42
GROUP BY o2.item_id ORDER BY COUNT(*) DESC LIMIT 30;
\end{lstlisting}
El plan usa el B+ de \texttt{ordenes.item\_id} para hallar a los compradores, un hash join externo, \texttt{GROUP BY} por hash y un top-N con heap que solo guarda los 30 mejores. Es la señal más fuerte y la más cara, unos 790~ms. En un producto sin historial (\emph{cold-start}) devuelve vacío.

\paragraph{Visual.} Compara la imagen del producto con las del catálogo, o una foto que sube el usuario. Hay tres modelos:
\begin{itemize}
  \item SIFT + BoVW con 256 palabras, el pipeline de la parte 4 \citep{lowe,sivic}: hasta 500 descriptores SIFT por imagen, vocabulario entrenado con el K-Means propio sobre 100\,000 descriptores del catálogo, histograma con TF-IDF e índice HNSW con coseno \citep{hnsw}. El motor extrae los descriptores de la consulta a partir de la ruta (unos 130~ms).
  \item SIFT + BoVW con 4\,096 palabras y TF sublineal. Como el motor carga un solo vocabulario a la vez, la aplicación calcula el vector reutilizando los módulos SIFT y TF-IDF del motor y lo busca con \texttt{SIMILAR\_TO(VECTOR(\ldots), k=30)}.
  \item CLIP ViT-B/32 \citep{clip}, una red preentrenada que da vectores de 512 dimensiones. Está solo para comparar contra SIFT.
\end{itemize}
Se agregaron los dos últimos porque con 256 palabras la búsqueda por foto fallaba en casos claros. Con la carátula de \emph{Pokémon Shield}, la de \emph{Pokémon Sword} quedaba en el puesto 273 de 1\,500: con tan pocas palabras todas las carátulas se parecen. Con la foto de un control DualShock negro, el blanco quedaba en el puesto 1\,352, porque SIFT describe la dirección del cambio de brillo y un borde negro sobre blanco produce descriptores casi opuestos a uno blanco sobre blanco. CLIP lo encuentra en el primer puesto.

\paragraph{Texto.} \texttt{MATCH(descripcion, \ldots) USING BM25 LIMIT 30} sobre el índice SPIMI de la parte 3 \citep{robertson}, usando como consulta las primeras 25 palabras del nombre, la marca, la categoría y la descripción del producto. El recorte mantiene el costo estable, en unos 30~ms.

\paragraph{Metadata.} Los productos mejor valorados de la misma categoría con un precio entre 0.5 y 1.8 veces el del producto consultado, completando con la misma categoría sin rango de precio si faltan candidatos. Usa el B+ de \texttt{categoria} y un top-N por rating (unos 31~ms).

\paragraph{Espacial.} La ubicación del usuario, elegida con un clic en el mapa, se usa de dos formas:
\begin{itemize}
  \item Como filtro: \texttt{distancia(ubicacion, POINT(\ldots)) < radio} o \texttt{intersecta(ubicacion, POLYGON(\ldots))} sobre el R-Tree \citep{guttman} devuelve las tiendas permitidas, que pasan a las demás señales como \texttt{sucursal\_id IN (\ldots)}. La distancia es Haversine \citep{sinnott}.
  \item Como señal: las 3 tiendas más cercanas (k-NN del R-Tree) dan a cada producto suyo un puntaje $1 / (1 + d / 5000\,\text{m})$. En el híbrido solo puntúa candidatos de las otras señales, para acercar el Top 10 sin agregar productos irrelevantes.
\end{itemize}

\subsection{Fusión híbrida}
Cada señal entrega sus candidatos con un puntaje en su propia escala (conteos, coseno, BM25 o rating). Se combinan con una de dos fórmulas, que el usuario elige:
\begin{itemize}
  \item Suma ponderada de puntajes normalizados, la opción por defecto:
  \[
  \text{puntaje}(p) = \sum_{m} w_m \, \frac{s_m(p)}{\max_q s_m(q)}
  \]
  La co-compra usa $\log(1 + \text{conteo})$ antes de normalizar, para que un producto con muchas co-compras no aplaste al resto. La señal espacial ya está entre 0 y 1 y no se normaliza.
  \item \emph{Reciprocal Rank Fusion} ponderada \citep{rrf}, la misma del RAG pero con un peso por señal:
  \[
  \text{RRF}(p) = \sum_{m} \frac{w_m}{60 + \text{rango}_m(p)}
  \]
\end{itemize}
Los pesos de las señales activas se renormalizan para sumar 1, y una señal con peso 0 no se consulta. Los empates se resuelven por mayor rating y luego por menor \texttt{id}. Los pesos por defecto se calibraron con los usuarios de validación (sección~\ref{sec:rec-eval}): co-compra 0.80, texto 0.15, visual 0.05 y metadata 0, más 0.10 de cercanía cuando hay ubicación. Cada recomendación devuelve cuánto aportó cada señal a su puntaje.

\subsection{Interfaz}
\begin{itemize}
  \item Catálogo paginado con filtros por categoría y precio, búsqueda por texto (BM25 o TF-IDF) y búsqueda por foto. En la búsqueda por foto, la imagen subida queda fija sobre los resultados y cada producto muestra su similitud coseno con ella.
  \item Al elegir un producto, un panel lateral muestra su ficha y el Top 10. Bajo cada recomendación, una franja de colores indica cuánto aportó cada señal.
  \item Un panel de preferencias permite cambiar la estrategia (híbrido o una señal sola), los pesos, la fórmula de fusión y el modelo de imagen sin cerrar el producto.
  \item Un mapa de Leaflet fija la ubicación del usuario, el radio y la zona, y marca las tiendas con recomendaciones y las más cercanas.
  \item Cada vista despliega el SQL enviado al motor, el plan ejecutado y los tiempos.
\end{itemize}
\imagen{imagenes/ecommerce/recomendacion.png}{Producto elegido y su Top 10 híbrido, con la franja de aporte de cada señal}{fig:rec-detalle}
\imagen{imagenes/ecommerce/busqueda_foto.png}{Búsqueda por foto: la imagen subida y los productos más parecidos con su similitud}{fig:rec-foto}

\subsection{Evaluación}\label{sec:rec-eval}
\subsubsection{Metodología}
\begin{itemize}
  \item Los usuarios con al menos 2 productos se repartieron al azar en entrenamiento (60\,\%), validación (20\,\%) y prueba (20\,\%), con una semilla fija.
  \item Por cada usuario de prueba (100 por corrida) se elige al azar uno de sus productos como consulta. Los objetivos son sus demás productos. La co-compra excluye las órdenes de los usuarios de prueba, para no usar la respuesta.
  \item El escenario \emph{cold-start} repite los mismos casos sin ninguna co-compra, como un producto nuevo.
  \item Los pesos del híbrido se eligieron solo con los 301 usuarios de validación, probando las 1\,771 combinaciones con paso 0.05 que suman 1. Se eligió la de mayor NDCG@10 promedio entre el escenario normal y el \emph{cold-start}. Optimizar solo el normal dejaba en 0 todas las señales de contenido.
  \item Métricas: Hit@10 (casos con al menos un objetivo en el Top 10), Recall@10, NDCG@10 \citep{ndcg}, MRR y Prec.~Cat@10, la fracción del Top 10 en la misma categoría que el producto consultado.
\end{itemize}

\subsubsection{Resultados}
\begin{table}[H]
\centering\small
\caption{Recomendación por señal y fusión, 100 usuarios de prueba}
\begin{tabular}{llccccc}
\toprule
Escenario & Modo & Hit@10 & Recall@10 & NDCG@10 & MRR & Prec. Cat@10 \\
\midrule
Normal & Co-compra & 0.60 & 0.083 & 0.134 & 0.324 & 0.304 \\
 & Visual (SIFT k=256) & 0.20 & 0.019 & 0.028 & 0.073 & 0.219 \\
 & Texto & 0.34 & 0.037 & 0.063 & 0.163 & 0.439 \\
 & Metadata & 0.15 & 0.014 & 0.017 & 0.035 & \textbf{1.000} \\
 & Híbrido RRF & \textbf{0.63} & 0.085 & \textbf{0.137} & \textbf{0.343} & 0.321 \\
 & Híbrido score & 0.61 & \textbf{0.086} & 0.135 & 0.323 & 0.323 \\
\midrule
Cold-start & Texto & \textbf{0.34} & \textbf{0.037} & \textbf{0.063} & \textbf{0.163} & 0.439 \\
 & Híbrido RRF & \textbf{0.34} & \textbf{0.037} & 0.060 & 0.145 & 0.440 \\
 & Híbrido score & \textbf{0.34} & \textbf{0.037} & 0.062 & 0.156 & 0.432 \\
\bottomrule
\end{tabular}
\end{table}
En \emph{cold-start} la co-compra no recomienda nada y las demás señales solas rinden igual que en el escenario normal. Con otra semilla de casos los resultados se mantienen: co-compra 0.61 de Hit@10 e híbrido score 0.63.

\begin{table}[H]
\centering\small
\caption{Modelos de imagen como señal sola, mismos casos}
\begin{tabular}{lccccc}
\toprule
Modelo & Hit@10 & Recall@10 & NDCG@10 & MRR & Prec. Cat@10 \\
\midrule
SIFT k=256 & 0.20 & 0.019 & 0.028 & 0.073 & 0.219 \\
SIFT k=4096 & 0.27 & 0.027 & 0.040 & 0.109 & \textbf{0.373} \\
CLIP & \textbf{0.38} & \textbf{0.044} & \textbf{0.067} & \textbf{0.171} & 0.288 \\
\bottomrule
\end{tabular}
\end{table}

\begin{table}[H]
\centering\small
\caption{Peso de la cercanía en el híbrido score, con el usuario en uno de 8 puntos de Lima}
\begin{tabular}{cccccc}
\toprule
Peso espacial & Hit@10 & NDCG@10 & MRR & Prec. Cat@10 & Distancia media del Top 10 \\
\midrule
0.00 & 0.61 & \textbf{0.135} & \textbf{0.323} & 0.323 & 7.86 km \\
0.10 & \textbf{0.62} & 0.130 & 0.297 & \textbf{0.328} & 6.77 km \\
0.25 & \textbf{0.62} & 0.121 & 0.279 & 0.318 & 5.63 km \\
0.50 & 0.58 & 0.107 & 0.247 & 0.295 & 4.68 km \\
\bottomrule
\end{tabular}
\end{table}

\begin{table}[H]
\centering\small
\caption{Tiempo en el motor de una recomendación híbrida con ubicación}
\begin{tabular}{llr}
\toprule
Consulta & Plan & Tiempo \\
\midrule
Co-compra (self-join) & B+ \texttt{item\_id}, hash join externo, top-N & $\sim$792 ms \\
Datos de los candidatos de co-compra & Sequential Scan & $\sim$127 ms \\
Visual & HNSW coseno & $\sim$133 ms \\
Texto & SPIMI BM25 & $\sim$30 ms \\
Metadata & B+ \texttt{categoria}, top-N & $\sim$31 ms \\
Radio y k-NN & R-Tree & $\sim$1 ms cada una \\
\bottomrule
\end{tabular}
\end{table}

\subsubsection{Análisis}
\begin{itemize}
  \item Con historial de compras, el híbrido iguala a la co-compra sola (Hit@10 entre 0.61 y 0.63 frente a 0.60 y 0.61) y sube algo la coherencia de categoría. Con 100 casos, diferencias de 0.02 están dentro del ruido.
  \item En \emph{cold-start} el híbrido rinde como el texto solo, la mejor señal en ese escenario. Con visual 0.05 y metadata 0, al faltar la co-compra casi solo queda el texto.
  \item Elegir los pesos es elegir una métrica. Unos pesos manuales con más metadata y visual (co-compra 0.55, texto 0.25, visual 0.08 y metadata 0.12) dan el mismo NDCG@10, pero más Hit@10 (0.66 en el escenario normal y 0.40 en \emph{cold-start}) y más coherencia de categoría (0.40 y 0.60, frente a 0.32 y 0.43). La calibración optimiza el NDCG sobre \emph{lo que el usuario también compró}, que favorece a la co-compra por diseño. Esos pesos manuales se eligieron antes de separar los usuarios al azar, así que su ventaja no es una medición limpia.
  \item Suma de puntajes y RRF quedan empatadas: cada una gana en una de las dos semillas, siempre dentro del ruido.
  \item Como señal sola, CLIP casi duplica el Hit@10 de SIFT con 256 palabras. SIFT con 4\,096 palabras tiene la mejor coherencia de categoría, porque reconoce el marco y el logo de la plataforma en la carátula. En el híbrido el modelo de imagen casi no cambia los resultados, porque su peso es 0.05; importa sobre todo en la búsqueda por foto.
  \item Un peso espacial de 0.10 acerca el Top 10 unos 1.1~km sin bajar el Hit@10. Desde 0.50 empeoran todas las métricas.
  \item La co-compra concentra cerca del 80\,\% del tiempo de una recomendación (unos 1\,000~ms en total). El R-Tree es prácticamente gratis.
\end{itemize}

\subsection{Alcance y limitaciones}
\begin{itemize}
  \item Un acierto significa que el usuario también compró el producto. La métrica mide solo de forma indirecta la similitud visual o de texto, y la coherencia de categoría la complementa.
  \item La co-compra no funciona con productos nuevos, y el self-join crece con la popularidad del producto. A mayor escala convendría precalcular los pares de co-compra.
  \item SIFT no encuentra el mismo objeto con otro color. CLIP sí, pero es un modelo preentrenado externo al minigestor.
  \item \texttt{WHERE id IN (\ldots)} sobre el catálogo se resuelve con un recorrido secuencial en vez del B+ de \texttt{id}. No se corrigió para no modificar el motor.
  \item La interfaz no registra compras, así que la co-compra no aprende en línea.
  \item La señal espacial en el híbrido solo reordena candidatos de las otras señales: no descubre productos nuevos.
\end{itemize}

% ==================== REFERENCIAS NUEVAS ====================
% Copiar dentro de \begin{thebibliography}{99} de main.tex, en orden alfabético.
% Las demás citas de esta sección (rrf, lowe, sivic, hnsw, robertson, guttman, sinnott) ya existen.
%
% \bibitem[Burke(2002)]{burke} Burke, R. (2002). Hybrid recommender systems: survey and experiments. \textit{User Modeling and User-Adapted Interaction}, 12(4), 331--370.
% \bibitem[Hou et~al.(2024)]{amazon2023} Hou, Y., Li, J., He, Z., Yan, A., Chen, X. y McAuley, J. (2024). Bridging language and items for retrieval and recommendation. \textit{arXiv:2403.03952}. Dataset: \url{https://amazon-reviews-2023.github.io/}
% \bibitem[Järvelin y Kekäläinen(2002)]{ndcg} Järvelin, K. y Kekäläinen, J. (2002). Cumulated gain-based evaluation of IR techniques. \textit{ACM Transactions on Information Systems}, 20(4), 422--446.
% \bibitem[Linden et~al.(2003)]{linden} Linden, G., Smith, B. y York, J. (2003). Amazon.com recommendations: item-to-item collaborative filtering. \textit{IEEE Internet Computing}, 7(1), 76--80.
% \bibitem[Radford et~al.(2021)]{clip} Radford, A., Kim, J. W., Hallacy, C., Ramesh, A., Goh, G., Agarwal, S., Sastry, G., Askell, A., Mishkin, P., Clark, J., Krueger, G. y Sutskever, I. (2021). Learning transferable visual models from natural language supervision. \textit{Proceedings of ICML}, PMLR 139, 8748--8763.
 justo después de la sección
% "Parte 5: Aplicación de inteligencia artificial" y antes de \begin{thebibliography}.
% Usa los macros del preámbulo de main.tex (\imagen, estilos TikZ caja/flecha, \citep).
% Las referencias nuevas están comentadas al final de este archivo: copiarlas a thebibliography.

\section{Parte 5: Sistema de recomendación e-commerce}
Además del RAG se implementó la opción B del enunciado: una tienda en línea de videojuegos que, para cada producto, recomienda un Top 10 combinando cinco señales \citep{burke}. Cada señal se resuelve con una estructura distinta del minigestor, así que la aplicación integra los cuatro tipos de datos del proyecto: relacional (órdenes, categoría y precio), espacial (ubicación de las tiendas), texto (descripciones) y multimedia (imágenes de los productos).

\begin{table}[H]
\centering\small
\caption{Señales del recomendador y estructura del minigestor que usa cada una}
\begin{tabularx}{\linewidth}{lYl}
\toprule
Señal & Pregunta que responde & Estructura \\
\midrule
Co-compra & ¿Qué compraron también quienes compraron este producto? & B+, self-join, \texttt{GROUP BY} \\
Texto & ¿Qué productos tienen una descripción parecida? & Índice invertido SPIMI, BM25 \\
Visual & ¿Qué productos se ven parecidos? & SIFT + BoVW, HNSW coseno \\
Metadata & ¿Qué productos de la misma categoría y precio están mejor valorados? & B+ sobre categoría \\
Espacial & ¿Qué productos hay en tiendas cerca de mí? & R-Tree: radio, k-NN y polígono \\
\bottomrule
\end{tabularx}
\end{table}

\subsection{Arquitectura}
La aplicación vive en la carpeta \texttt{ecommerce/} del repositorio y es independiente del motor. Se comunica con él solo mediante \texttt{POST /api/query}, igual que el RAG, y no modifica el código del minigestor.

\begin{figure}[H]
\centering
\resizebox{\linewidth}{!}{%
\begin{tikzpicture}[node distance=12mm and 14mm]
  \node[caja, minimum width=3cm] (nav) {\textbf{Tienda web}\\HTML + JS + Leaflet};
  \node[caja, right=of nav, minimum width=3.4cm] (srv) {\textbf{Servidor de la tienda}\\FastAPI, puerto 8300};
  \node[caja, right=of srv, minimum width=3.4cm] (rec) {\textbf{Recomendador}\\5 señales $\cdot$ fusión};
  \node[caja, below=of rec, minimum width=3.4cm] (vis) {\textbf{Vectores de imagen}\\SIFT k=4096 $\cdot$ CLIP\\local, CPU};
  \node[caja, below=of srv, minimum width=3.4cm, fill=gray!10] (motor) {\textbf{Minigestor BD2}\\puerto 8000\\B+ $\cdot$ SPIMI $\cdot$ HNSW $\cdot$ R-Tree};
  \draw[flecha] (nav) -- node[above, font=\scriptsize]{producto o foto} (srv);
  \draw[flecha] (srv) -- (rec);
  \draw[flecha] (rec) -- node[right, font=\scriptsize]{imagen} (vis);
  \draw[flecha] (rec) -- node[above left, font=\scriptsize]{SQL por HTTP} (motor);
  \draw[flecha, dashed] (rec.north) -- ++(0,0.6) -| node[above, pos=0.25, font=\scriptsize]{Top 10 con el aporte de cada señal} (nav.north);
\end{tikzpicture}}
\caption{Arquitectura del recomendador: la aplicación solo usa la API SQL del minigestor}\label{fig:rec}
\end{figure}

\begin{itemize}
  \item \texttt{recomendador.py} tiene una función por señal y la fusión. No guarda estado: cada recomendación consulta al motor.
  \item \texttt{cliente\_bd.py} es la única puerta hacia el motor. Arma los literales SQL de forma segura y anota el SQL, el plan y los tiempos de cada consulta, que la interfaz muestra.
  \item \texttt{visuales.py} calcula en la aplicación los vectores de los dos modelos de imagen que el motor no tiene cargados (sección~\ref{sec:rec-senales}).
\end{itemize}

\subsection{Dataset}
\begin{itemize}
  \item Se usó Amazon Reviews 2023 \citep{amazon2023}, categoría \texttt{Video\_Games}. Las interacciones salen del archivo \emph{5-core} (usuarios y productos con al menos 5 reseñas) y se tomaron los cerca de 4\,500 productos más reseñados como candidatos.
  \item El catálogo final tiene 1\,500 productos con imagen, precio (puede faltar), rating, marca y categoría. La categoría usa el segundo y tercer nivel del árbol de Amazon, por ejemplo \emph{Nintendo Switch > Games}, así que queda separada por plataforma.
  \item La descripción une título, marca, categoría, características y descripción del producto en un solo texto.
  \item Cada reseña se toma como una compra. Se generó una orden por par (usuario, producto), solo para usuarios con al menos 2 productos del catálogo: resultaron 24\,999 órdenes.
  \item Las 8 tiendas son puntos reales de Lima (Miraflores, San Isidro, Surco, San Borja, Barranco, Jesús María, San Miguel y Callao). Cada producto se asigna a una tienda de forma determinista con \texttt{crc32(codigo) mod 8}.
\end{itemize}

\subsection{Modelo de datos}
\begin{lstlisting}[language=SQL]
CREATE TABLE sucursales (id INT PRIMARY KEY, nombre VARCHAR(100),
  distrito VARCHAR(60), ubicacion POINT) USING BPLUS;
CREATE TABLE catalogo (id INT PRIMARY KEY, codigo VARCHAR(20),
  nombre VARCHAR(200), marca VARCHAR(80), categoria VARCHAR(80),
  precio FLOAT, rating FLOAT, sucursal_id INT, imagen VARCHAR(255),
  descripcion TEXT, emb VECTOR(256)) USING BPLUS;
CREATE TABLE ordenes (id INT PRIMARY KEY, usuario_id INT, item_id INT,
  fecha DATE) USING BPLUS;
CREATE INDEX idx_cat_categoria ON catalogo (categoria) USING BPLUS;
CREATE INDEX idx_cat_sucursal ON catalogo (sucursal_id) USING BPLUS;
CREATE INDEX ft_catalogo ON catalogo (descripcion) USING FULLTEXT;
CREATE INDEX hnsw_catalogo ON catalogo (emb) USING HNSW WITH (METRIC=cosine);
CREATE INDEX idx_ord_usuario ON ordenes (usuario_id) USING BPLUS;
CREATE INDEX idx_ord_item ON ordenes (item_id) USING BPLUS;
\end{lstlisting}
\begin{itemize}
  \item La columna \texttt{ubicacion} es \texttt{POINT}, así que el motor le crea su R-Tree automáticamente (parte 2).
  \item El vector \texttt{emb} lo calcula el propio motor a partir de la ruta de la imagen, con el pipeline SIFT de la parte 4.
  \item Los otros dos modelos de imagen viven en tablas propias, \texttt{vis\_sift4096} con \texttt{VECTOR(4096)} y \texttt{vis\_clip} con \texttt{VECTOR(512)}, cada una con su índice HNSW. Son opcionales.
  \item La ingesta inserta por lotes de 40 filas, crea los índices al final y corre \texttt{ANALYZE} en cada tabla para que el planificador tenga estadísticas.
\end{itemize}

\subsection{Señales de recomendación}\label{sec:rec-senales}
Todas las señales reciben los mismos filtros de categoría, precio y tiendas permitidas como condiciones extra del \texttt{WHERE}, excluyen al producto consultado y devuelven hasta 30 candidatos.

\paragraph{Co-compra.} Es filtrado colaborativo de producto a producto \citep{linden}: \emph{quienes compraron X también compraron Y}.
\begin{lstlisting}[language=SQL]
SELECT o2.item_id, COUNT(*) FROM ordenes AS o1
JOIN ordenes AS o2 ON o1.usuario_id = o2.usuario_id
WHERE o1.item_id = 42 AND o2.item_id != 42
GROUP BY o2.item_id ORDER BY COUNT(*) DESC LIMIT 30;
\end{lstlisting}
El plan usa el B+ de \texttt{ordenes.item\_id} para hallar a los compradores, un hash join externo, \texttt{GROUP BY} por hash y un top-N con heap que solo guarda los 30 mejores. Es la señal más fuerte y la más cara, unos 790~ms. En un producto sin historial (\emph{cold-start}) devuelve vacío.

\paragraph{Visual.} Compara la imagen del producto con las del catálogo, o una foto que sube el usuario. Hay tres modelos:
\begin{itemize}
  \item SIFT + BoVW con 256 palabras, el pipeline de la parte 4 \citep{lowe,sivic}: hasta 500 descriptores SIFT por imagen, vocabulario entrenado con el K-Means propio sobre 100\,000 descriptores del catálogo, histograma con TF-IDF e índice HNSW con coseno \citep{hnsw}. El motor extrae los descriptores de la consulta a partir de la ruta (unos 130~ms).
  \item SIFT + BoVW con 4\,096 palabras y TF sublineal. Como el motor carga un solo vocabulario a la vez, la aplicación calcula el vector reutilizando los módulos SIFT y TF-IDF del motor y lo busca con \texttt{SIMILAR\_TO(VECTOR(\ldots), k=30)}.
  \item CLIP ViT-B/32 \citep{clip}, una red preentrenada que da vectores de 512 dimensiones. Está solo para comparar contra SIFT.
\end{itemize}
Se agregaron los dos últimos porque con 256 palabras la búsqueda por foto fallaba en casos claros. Con la carátula de \emph{Pokémon Shield}, la de \emph{Pokémon Sword} quedaba en el puesto 273 de 1\,500: con tan pocas palabras todas las carátulas se parecen. Con la foto de un control DualShock negro, el blanco quedaba en el puesto 1\,352, porque SIFT describe la dirección del cambio de brillo y un borde negro sobre blanco produce descriptores casi opuestos a uno blanco sobre blanco. CLIP lo encuentra en el primer puesto.

\paragraph{Texto.} \texttt{MATCH(descripcion, \ldots) USING BM25 LIMIT 30} sobre el índice SPIMI de la parte 3 \citep{robertson}, usando como consulta las primeras 25 palabras del nombre, la marca, la categoría y la descripción del producto. El recorte mantiene el costo estable, en unos 30~ms.

\paragraph{Metadata.} Los productos mejor valorados de la misma categoría con un precio entre 0.5 y 1.8 veces el del producto consultado, completando con la misma categoría sin rango de precio si faltan candidatos. Usa el B+ de \texttt{categoria} y un top-N por rating (unos 31~ms).

\paragraph{Espacial.} La ubicación del usuario, elegida con un clic en el mapa, se usa de dos formas:
\begin{itemize}
  \item Como filtro: \texttt{distancia(ubicacion, POINT(\ldots)) < radio} o \texttt{intersecta(ubicacion, POLYGON(\ldots))} sobre el R-Tree \citep{guttman} devuelve las tiendas permitidas, que pasan a las demás señales como \texttt{sucursal\_id IN (\ldots)}. La distancia es Haversine \citep{sinnott}.
  \item Como señal: las 3 tiendas más cercanas (k-NN del R-Tree) dan a cada producto suyo un puntaje $1 / (1 + d / 5000\,\text{m})$. En el híbrido solo puntúa candidatos de las otras señales, para acercar el Top 10 sin agregar productos irrelevantes.
\end{itemize}

\subsection{Fusión híbrida}
Cada señal entrega sus candidatos con un puntaje en su propia escala (conteos, coseno, BM25 o rating). Se combinan con una de dos fórmulas, que el usuario elige:
\begin{itemize}
  \item Suma ponderada de puntajes normalizados, la opción por defecto:
  \[
  \text{puntaje}(p) = \sum_{m} w_m \, \frac{s_m(p)}{\max_q s_m(q)}
  \]
  La co-compra usa $\log(1 + \text{conteo})$ antes de normalizar, para que un producto con muchas co-compras no aplaste al resto. La señal espacial ya está entre 0 y 1 y no se normaliza.
  \item \emph{Reciprocal Rank Fusion} ponderada \citep{rrf}, la misma del RAG pero con un peso por señal:
  \[
  \text{RRF}(p) = \sum_{m} \frac{w_m}{60 + \text{rango}_m(p)}
  \]
\end{itemize}
Los pesos de las señales activas se renormalizan para sumar 1, y una señal con peso 0 no se consulta. Los empates se resuelven por mayor rating y luego por menor \texttt{id}. Los pesos por defecto se calibraron con los usuarios de validación (sección~\ref{sec:rec-eval}): co-compra 0.80, texto 0.15, visual 0.05 y metadata 0, más 0.10 de cercanía cuando hay ubicación. Cada recomendación devuelve cuánto aportó cada señal a su puntaje.

\subsection{Interfaz}
\begin{itemize}
  \item Catálogo paginado con filtros por categoría y precio, búsqueda por texto (BM25 o TF-IDF) y búsqueda por foto. En la búsqueda por foto, la imagen subida queda fija sobre los resultados y cada producto muestra su similitud coseno con ella.
  \item Al elegir un producto, un panel lateral muestra su ficha y el Top 10. Bajo cada recomendación, una franja de colores indica cuánto aportó cada señal.
  \item Un panel de preferencias permite cambiar la estrategia (híbrido o una señal sola), los pesos, la fórmula de fusión y el modelo de imagen sin cerrar el producto.
  \item Un mapa de Leaflet fija la ubicación del usuario, el radio y la zona, y marca las tiendas con recomendaciones y las más cercanas.
  \item Cada vista despliega el SQL enviado al motor, el plan ejecutado y los tiempos.
\end{itemize}
\imagen{imagenes/ecommerce/recomendacion.png}{Producto elegido y su Top 10 híbrido, con la franja de aporte de cada señal}{fig:rec-detalle}
\imagen{imagenes/ecommerce/busqueda_foto.png}{Búsqueda por foto: la imagen subida y los productos más parecidos con su similitud}{fig:rec-foto}

\subsection{Evaluación}\label{sec:rec-eval}
\subsubsection{Metodología}
\begin{itemize}
  \item Los usuarios con al menos 2 productos se repartieron al azar en entrenamiento (60\,\%), validación (20\,\%) y prueba (20\,\%), con una semilla fija.
  \item Por cada usuario de prueba (100 por corrida) se elige al azar uno de sus productos como consulta. Los objetivos son sus demás productos. La co-compra excluye las órdenes de los usuarios de prueba, para no usar la respuesta.
  \item El escenario \emph{cold-start} repite los mismos casos sin ninguna co-compra, como un producto nuevo.
  \item Los pesos del híbrido se eligieron solo con los 301 usuarios de validación, probando las 1\,771 combinaciones con paso 0.05 que suman 1. Se eligió la de mayor NDCG@10 promedio entre el escenario normal y el \emph{cold-start}. Optimizar solo el normal dejaba en 0 todas las señales de contenido.
  \item Métricas: Hit@10 (casos con al menos un objetivo en el Top 10), Recall@10, NDCG@10 \citep{ndcg}, MRR y Prec.~Cat@10, la fracción del Top 10 en la misma categoría que el producto consultado.
\end{itemize}

\subsubsection{Resultados}
\begin{table}[H]
\centering\small
\caption{Recomendación por señal y fusión, 100 usuarios de prueba}
\begin{tabular}{llccccc}
\toprule
Escenario & Modo & Hit@10 & Recall@10 & NDCG@10 & MRR & Prec. Cat@10 \\
\midrule
Normal & Co-compra & 0.60 & 0.083 & 0.134 & 0.324 & 0.304 \\
 & Visual (SIFT k=256) & 0.20 & 0.019 & 0.028 & 0.073 & 0.219 \\
 & Texto & 0.34 & 0.037 & 0.063 & 0.163 & 0.439 \\
 & Metadata & 0.15 & 0.014 & 0.017 & 0.035 & \textbf{1.000} \\
 & Híbrido RRF & \textbf{0.63} & 0.085 & \textbf{0.137} & \textbf{0.343} & 0.321 \\
 & Híbrido score & 0.61 & \textbf{0.086} & 0.135 & 0.323 & 0.323 \\
\midrule
Cold-start & Texto & \textbf{0.34} & \textbf{0.037} & \textbf{0.063} & \textbf{0.163} & 0.439 \\
 & Híbrido RRF & \textbf{0.34} & \textbf{0.037} & 0.060 & 0.145 & 0.440 \\
 & Híbrido score & \textbf{0.34} & \textbf{0.037} & 0.062 & 0.156 & 0.432 \\
\bottomrule
\end{tabular}
\end{table}
En \emph{cold-start} la co-compra no recomienda nada y las demás señales solas rinden igual que en el escenario normal. Con otra semilla de casos los resultados se mantienen: co-compra 0.61 de Hit@10 e híbrido score 0.63.

\begin{table}[H]
\centering\small
\caption{Modelos de imagen como señal sola, mismos casos}
\begin{tabular}{lccccc}
\toprule
Modelo & Hit@10 & Recall@10 & NDCG@10 & MRR & Prec. Cat@10 \\
\midrule
SIFT k=256 & 0.20 & 0.019 & 0.028 & 0.073 & 0.219 \\
SIFT k=4096 & 0.27 & 0.027 & 0.040 & 0.109 & \textbf{0.373} \\
CLIP & \textbf{0.38} & \textbf{0.044} & \textbf{0.067} & \textbf{0.171} & 0.288 \\
\bottomrule
\end{tabular}
\end{table}

\begin{table}[H]
\centering\small
\caption{Peso de la cercanía en el híbrido score, con el usuario en uno de 8 puntos de Lima}
\begin{tabular}{cccccc}
\toprule
Peso espacial & Hit@10 & NDCG@10 & MRR & Prec. Cat@10 & Distancia media del Top 10 \\
\midrule
0.00 & 0.61 & \textbf{0.135} & \textbf{0.323} & 0.323 & 7.86 km \\
0.10 & \textbf{0.62} & 0.130 & 0.297 & \textbf{0.328} & 6.77 km \\
0.25 & \textbf{0.62} & 0.121 & 0.279 & 0.318 & 5.63 km \\
0.50 & 0.58 & 0.107 & 0.247 & 0.295 & 4.68 km \\
\bottomrule
\end{tabular}
\end{table}

\begin{table}[H]
\centering\small
\caption{Tiempo en el motor de una recomendación híbrida con ubicación}
\begin{tabular}{llr}
\toprule
Consulta & Plan & Tiempo \\
\midrule
Co-compra (self-join) & B+ \texttt{item\_id}, hash join externo, top-N & $\sim$792 ms \\
Datos de los candidatos de co-compra & Sequential Scan & $\sim$127 ms \\
Visual & HNSW coseno & $\sim$133 ms \\
Texto & SPIMI BM25 & $\sim$30 ms \\
Metadata & B+ \texttt{categoria}, top-N & $\sim$31 ms \\
Radio y k-NN & R-Tree & $\sim$1 ms cada una \\
\bottomrule
\end{tabular}
\end{table}

\subsubsection{Análisis}
\begin{itemize}
  \item Con historial de compras, el híbrido iguala a la co-compra sola (Hit@10 entre 0.61 y 0.63 frente a 0.60 y 0.61) y sube algo la coherencia de categoría. Con 100 casos, diferencias de 0.02 están dentro del ruido.
  \item En \emph{cold-start} el híbrido rinde como el texto solo, la mejor señal en ese escenario. Con visual 0.05 y metadata 0, al faltar la co-compra casi solo queda el texto.
  \item Elegir los pesos es elegir una métrica. Unos pesos manuales con más metadata y visual (co-compra 0.55, texto 0.25, visual 0.08 y metadata 0.12) dan el mismo NDCG@10, pero más Hit@10 (0.66 en el escenario normal y 0.40 en \emph{cold-start}) y más coherencia de categoría (0.40 y 0.60, frente a 0.32 y 0.43). La calibración optimiza el NDCG sobre \emph{lo que el usuario también compró}, que favorece a la co-compra por diseño. Esos pesos manuales se eligieron antes de separar los usuarios al azar, así que su ventaja no es una medición limpia.
  \item Suma de puntajes y RRF quedan empatadas: cada una gana en una de las dos semillas, siempre dentro del ruido.
  \item Como señal sola, CLIP casi duplica el Hit@10 de SIFT con 256 palabras. SIFT con 4\,096 palabras tiene la mejor coherencia de categoría, porque reconoce el marco y el logo de la plataforma en la carátula. En el híbrido el modelo de imagen casi no cambia los resultados, porque su peso es 0.05; importa sobre todo en la búsqueda por foto.
  \item Un peso espacial de 0.10 acerca el Top 10 unos 1.1~km sin bajar el Hit@10. Desde 0.50 empeoran todas las métricas.
  \item La co-compra concentra cerca del 80\,\% del tiempo de una recomendación (unos 1\,000~ms en total). El R-Tree es prácticamente gratis.
\end{itemize}

\subsection{Alcance y limitaciones}
\begin{itemize}
  \item Un acierto significa que el usuario también compró el producto. La métrica mide solo de forma indirecta la similitud visual o de texto, y la coherencia de categoría la complementa.
  \item La co-compra no funciona con productos nuevos, y el self-join crece con la popularidad del producto. A mayor escala convendría precalcular los pares de co-compra.
  \item SIFT no encuentra el mismo objeto con otro color. CLIP sí, pero es un modelo preentrenado externo al minigestor.
  \item \texttt{WHERE id IN (\ldots)} sobre el catálogo se resuelve con un recorrido secuencial en vez del B+ de \texttt{id}. No se corrigió para no modificar el motor.
  \item La interfaz no registra compras, así que la co-compra no aprende en línea.
  \item La señal espacial en el híbrido solo reordena candidatos de las otras señales: no descubre productos nuevos.
\end{itemize}

% ==================== REFERENCIAS NUEVAS ====================
% Copiar dentro de \begin{thebibliography}{99} de main.tex, en orden alfabético.
% Las demás citas de esta sección (rrf, lowe, sivic, hnsw, robertson, guttman, sinnott) ya existen.
%
% \bibitem[Burke(2002)]{burke} Burke, R. (2002). Hybrid recommender systems: survey and experiments. \textit{User Modeling and User-Adapted Interaction}, 12(4), 331--370.
% \bibitem[Hou et~al.(2024)]{amazon2023} Hou, Y., Li, J., He, Z., Yan, A., Chen, X. y McAuley, J. (2024). Bridging language and items for retrieval and recommendation. \textit{arXiv:2403.03952}. Dataset: \url{https://amazon-reviews-2023.github.io/}
% \bibitem[Järvelin y Kekäläinen(2002)]{ndcg} Järvelin, K. y Kekäläinen, J. (2002). Cumulated gain-based evaluation of IR techniques. \textit{ACM Transactions on Information Systems}, 20(4), 422--446.
% \bibitem[Linden et~al.(2003)]{linden} Linden, G., Smith, B. y York, J. (2003). Amazon.com recommendations: item-to-item collaborative filtering. \textit{IEEE Internet Computing}, 7(1), 76--80.
% \bibitem[Radford et~al.(2021)]{clip} Radford, A., Kim, J. W., Hallacy, C., Ramesh, A., Goh, G., Agarwal, S., Sastry, G., Askell, A., Mishkin, P., Clark, J., Krueger, G. y Sutskever, I. (2021). Learning transferable visual models from natural language supervision. \textit{Proceedings of ICML}, PMLR 139, 8748--8763.
